% Generated draft: requires mathematical and source review.
\documentclass[11pt]{article}
\usepackage[margin=25mm]{geometry}
\usepackage{fontspec}
\setmainfont{Latin Modern Roman}
\usepackage{amsmath,amssymb,mathtools,unicode-math}
\setmathfont{Latin Modern Math}
\usepackage{enumitem,needspace,xurl,hyperref}
\setlength{\parindent}{0pt}
\setlength{\parskip}{5pt}
\setlength{\emergencystretch}{3em}
\newcommand{\N}{\mathbb{N}}
\newcommand{\Z}{\mathbb{Z}}
\newcommand{\Q}{\mathbb{Q}}
\newcommand{\R}{\mathbb{R}}
\newcommand{\C}{\mathbb{C}}
\newcommand{\F}{\mathbb{F}}
\newcommand{\A}{\mathbb{A}}
\newcommand{\PP}{\mathbb{P}}
\newcommand{\E}{\mathbb{E}}
\newcommand{\eps}{\varepsilon}
\newcommand{\dd}{\,\mathrm{d}}
\begin{document}
\section*{UnsolvedMath ID: 20003184}
\textbf{Source problem code:} catalogue entry
\subsection{Definitions and mathematical statement}
Fix a closed Riemann surface $B$ of genus $g\geq3$ and a closed oriented surface $F_h$ of genus $h\geq2$. Dieter Kotschick asks which conjugacy classes
\[
\rho:\pi_1(B)\to\operatorname{MCG}(F_h)
\]
are monodromies of smooth holomorphic genus-$h$ fibrations $F_h\to X\to B$, meaning proper holomorphic submersions over this fixed complex base. The source records finiteness, up to the intended isomorphism over $B$, of the nonisotrivial such fibrations. Its Comment 3.2 notes that $X$ is K\"ahler, giving Hodge-theoretic restrictions such as even first Betti number $b_1(X)=2h^{1,0}(X)$, but says these are not sufficient for an if-and-only-if characterization. The extracted $h_1(X)$ is interpreted as this Betti-number assertion, not a claim that a Hodge number must be even.

Let $\widetilde B$ be the universal cover and $\mathcal T_h$ Teichm\"uller space of the fiber. The aid's exact geometric reformulation is
\[
\rho\text{ is holomorphically realizable over }B
\quad\Longleftrightarrow\quad
\exists\text{ holomorphic }u:\widetilde B\to\mathcal T_h,
\quad u(\gamma z)=\rho(\gamma)u(z).
\]
Equivalently the associated holomorphic Teichm\"uller bundle
\[
(\widetilde B\times\mathcal T_h)/\pi_1(B)\longrightarrow B
\]
has a holomorphic section, with diagonal action $\gamma(z,x)=(\gamma z,\rho(\gamma)x)$. This keeps the given complex structure on $B$, rather than allowing an arbitrary base deformation.

The aid reports that finite image is exactly the isotrivial case, that finite-image representations are realizable by the appropriate fixed-fiber construction, and that infinite reducible and virtually cyclic images cannot give the required nonisotrivial fibrations. Full image hypotheses and references for these exclusions are not supplied.

For the class called ``sufficiently large'' in the aid, assume the stated existence and uniqueness theorem supplies an interior Weil--Petersson harmonic equivariant map $u_\rho$. With compatible domain and target metrics, define its antiholomorphic defect by
\[
D_B(\rho)=\int_B\|\overline\partial u_\rho\|^2\,\dd A\geq0.
\]
The integrand descends by equivariance. The asserted analytic criterion is that $\rho$ is holomorphically realizable exactly when $D_B(\rho)=0$. The exact large-image hypotheses, harmonic-map regularity, interior image and uniqueness conditions must be checked in a definitive theorem; they are not inferred for arbitrary representations into an incomplete target.

For a connected unramified degree-$m$ cover $p:B'\to B$ with the induced complex structure and $\rho'=\rho|_{\pi_1(B')}$, the aid asserts in this large-image setting
\[
D_{B'}(\rho')=mD_B(\rho).
\]
Consequently realizability over $B$ is equivalent to realizability of the restricted representation over that $B'$. The aid states the same equivalence in the finite-image isotrivial case. This does not assert a branched-cover theorem, permit an unrelated complex structure on $B'$, or turn the reformulation into a finite algebraic test from the original representation alone.
\subsection{Short English statement}
Characterize holomorphic surface-bundle monodromy over the fixed Riemann surface. The supplied aid reformulates realization as a holomorphic Teichm\"uller section and, under its large-image harmonic-map hypotheses, as vanishing antiholomorphic energy, with an exact finite-cover scaling law.
\subsection{Sources}
\begin{itemize}
\item Local UnsolvedMath registry, permanent ID $20003184$, source code \texttt{catalogue entry}; canonical \texttt{aim-topology-notes.json}, \texttt{notes[271]}.
\item AIM workshop \emph{Representations of surface groups}, March 19--23, 2007, Question 3.1 and Comment 3.2, both attributed to Dieter Kotschick: \url{https://aimath.org/WWN/surfacegroups/surfacegroups.pdf}. Original remarks and literature field are empty; statement status is \texttt{reconstructed\_unverified}.
\item The aid invokes Teichm\"uller equivariant maps, finite-image realization and Weil--Petersson harmonic-map theory without authoritative complete citations or a precise sufficiently-large-image definition.
\item Research attempt 1: \texttt{reduction}, status \texttt{partially\_solved}, classification \texttt{PARTIAL-PROGRESS}, rigor \texttt{proved}, novelty confidence low. Its section criterion and finite-cover defect theorem are machine-generated aid, not peer reviewed, and require independent verification before citation.
\end{itemize}
\end{document}
